\subsection{Analytical and empirical goals}
In this paper, I draw out the empirical implications of this set of starting ideas for our understanding of the \emph{social uses} of aesthetic consumption in interaction.  My analytical point of departure is the proposition that one of the primary ``uses'' of aesthetic consumption are social \citep{dimaggio87};  even if the ``gratifications'' that it provides can sometimes be personal and private  \citep{denora00, hennion01, holbrook_hirschman82}.  This implies that the everyday use of cultural knowledge and experiences gained through aesthetic consumption become a strategic---although naturally not the only---resource for the initiation social interaction in conversation.  This happens in both relatively formal \citep{long03} and informal \citep{fine77} settings.   I propose that inequalities in cultural access can be more profitably understood as reflecting inequalities in the ability and opportunity to engage in certain forms of informal interaction that Simmel \citeyearpar{simmel49} referred to by the name of ``sociability.''  However, rather than thinking of sociability as simply a ``form'' of association, I focus on sociability as a status-group marker. I argue that one of the primary ways in which cultural advantage translates into social advantage is via the differential distribution of dispositions to generate sociable intercourse keyed to the mobilization of discourse associated with aesthetic consumption experiences.

